\documentclass[11pt,a4paper]{article}
\usepackage{fontspec,amsmath,amssymb,amsthm,microtype}
\setmainfont{lmroman10-regular.otf}[BoldFont=lmroman10-bold.otf,ItalicFont=lmroman10-italic.otf,BoldItalicFont=lmroman10-bolditalic.otf]
\setmonofont{lmmono10-regular.otf}[Scale=MatchLowercase]
\usepackage[margin=25mm]{geometry}
\usepackage[unicode,hidelinks]{hyperref}
\setlength{\emergencystretch}{3em}
\setlength{\parskip}{0.4em}
\setlength{\parindent}{0pt}
\allowdisplaybreaks[2]
\raggedbottom
\title{\texorpdfstring{One canonical global kinetic reference behind all coupling segments}{One canonical global kinetic reference behind all coupling segments}}
\author{LangevinLean mathematical exposition}
\date{}
\begin{document}\maketitle
\noindent Source module: \texttt{Kinetic\allowbreak{}Langevin/\allowbreak{}Canonical\allowbreak{}Global\allowbreak{}Segments.\allowbreak{}lean}. The accompanying HTML pairs this mathematical exposition with the exact Lean statements and records the scope of the formalization.


The exact reference segments are constructed on one canonical Brownian source. This module identifies their increment input with the increments of one full L\ensuremath{{}^{2}} Brownian input and proves that every stitched reference state and intermediate segment point belongs to the actual full kinetic solution.


\section{\texorpdfstring{Construct the full canonical input and identify its increments}{Construct the full canonical input and identify its increments}}


The full input is the scalar canonical Brownian value assembled spatially and pulled onto the actual initial-coupling extension with amplitude \(\sqrt{2\gamma}\). It is continuous in L\ensuremath{{}^{2}} because each assembly and pullback is continuous linear. The Brownian value at zero vanishes: Lean verifies that its deterministic time indicator has empty support, then uses linearity of the isonormal map. For nonnegative local time, the constructed segment input equals the full input difference at \(nh+t\) and \(nh\). The proof checks the nonnegative real-to-NNReal casts and uses linearity of spatial assembly and source pullback.



{\raggedright\noindent Lean declarations: \texttt{global\allowbreak{}Input}, \texttt{continuous\_\allowbreak{}global\allowbreak{}Input}, \texttt{global\allowbreak{}Input\_\allowbreak{}zero}, \texttt{input\_\allowbreak{}eq\_\allowbreak{}global\allowbreak{}Input\_\allowbreak{}increment}.\par}

\section{\texorpdfstring{Construct the actual full reference solution}{Construct the actual full reference solution}}


The global reference is the selected kinetic solution with the genuine lifted gradient, the canonical reference initial phase and the full canonical Brownian input. Its specification proves continuity, its exact starting phase, and the actual kinetic integral equation over the whole finite horizon. There is no stationary-law requirement in its definition.



{\raggedright\noindent Lean declarations: \texttt{global\allowbreak{}Reference\allowbreak{}Flow}, \texttt{global\allowbreak{}Reference\allowbreak{}Flow\_\allowbreak{}spec}.\par}

\section{\texorpdfstring{Identify every stitched reference point by uniqueness}{Identify every stitched reference point by uniqueness}}


The generic stitching theorem now applies because the local inputs are proved to be actual increments of the full input. Consequently \(R_n=Y_{nh}\), and \(Z_n(t)=Y_{nh+t}\) throughout each step whose endpoint lies in the global horizon. Lean instantiates the existing exact uniqueness results with the actual L\ensuremath{{}^{2}} gradient and checks the horizon arithmetic. Thus a stationary law for the full constructed kinetic solution supplies the law of all reference states used by the quantitative KLMC proof.



{\raggedright\noindent Lean declarations: \texttt{reference\allowbreak{}Chain\_\allowbreak{}eq\_\allowbreak{}global\allowbreak{}Reference\allowbreak{}Flow}, \texttt{reference\allowbreak{}Segment\_\allowbreak{}eq\_\allowbreak{}global\allowbreak{}Reference\allowbreak{}Flow}.\par}

\end{document}
